\documentclass[11pt]{article}

% --- Encoding & fonts (matches series style) ---
\usepackage[utf8]{inputenc}
\usepackage[T1]{fontenc}
\usepackage{lmodern}
\usepackage[protrusion=true,expansion=false]{microtype}
\setlength{\emergencystretch}{3em}

% --- Page geometry ---
\usepackage[margin=1in, top=1.1in, bottom=1.1in]{geometry}

% --- Packages ---
\usepackage{amsmath,amssymb,amsthm}
\usepackage{booktabs}
\usepackage{siunitx}
\usepackage{graphicx}
\usepackage[colorlinks=true,
            linkcolor=blue!70!black,
            citecolor=blue!70!black,
            urlcolor=blue!70!black]{hyperref}
\usepackage{xcolor}
\usepackage{enumitem}
\usepackage[numbers,sort&compress]{natbib}
\usepackage{caption}
\captionsetup{font=small,labelfont=bf}
\usepackage{tikz}
\usetikzlibrary{shapes.geometric,arrows.meta,positioning,fit,backgrounds}
\usepackage{multirow}
\usepackage{array}
\usepackage{algorithm}
\usepackage{algpseudocode}

% --- Section formatting ---
\usepackage{titlesec}
\titleformat{\section}{\large\bfseries}{\thesection}{1em}{}
\titleformat{\subsection}{\normalsize\bfseries}{\thesubsection}{1em}{}

% --- Paragraph spacing ---
\setlength{\parskip}{0.5em}
\setlength{\parindent}{0pt}

% --- Theorem environments ---
\newtheorem{theorem}{Theorem}[section]
\newtheorem{lemma}[theorem]{Lemma}
\newtheorem{corollary}[theorem]{Corollary}
\newtheorem{proposition}[theorem]{Proposition}
\newtheorem{definition}{Definition}[section]
\newtheorem{remark}{Remark}[section]

% --- Notation (inherits P7--P9 commands) ---
\newcommand{\Azero}{\mathcal{A}_0}
\newcommand{\Dhat}{\widehat{D}}
\newcommand{\HALT}{\texttt{HALT}}
\newcommand{\RESUME}{\texttt{RESUME}}
\newcommand{\DENY}{\texttt{DENY}}
\newcommand{\ADMIT}{\texttt{ADMIT}}
\newcommand{\ESCALATE}{\texttt{ESCALATE}}
\newcommand{\Pset}{\mathcal{P}}
\newcommand{\Es}{E_s}
\newcommand{\Dh}{D_h}
\newcommand{\sigh}{\sigma_h}
\newcommand{\Hi}{H_i}
\newcommand{\ski}{sk_i}
\newcommand{\pki}{pk_i}
\newcommand{\APB}{\mathrm{APB}}
% --- P9 notation ---
\newcommand{\Proxy}{\Pi}
\newcommand{\Dnet}{\Delta_{\mathrm{net}}}
\newcommand{\Dverify}{\Delta_{\mathrm{verify}}}
\newcommand{\Dpolicy}{\Delta_{\mathrm{policy}}}
% --- P10-specific ---
\newcommand{\Eescrow}{E_{\mathrm{escrow}}}
\newcommand{\thalt}{t_{\mathrm{halt}}}
\newcommand{\tsign}{t_{\mathrm{sign}}}
\newcommand{\tapb}{\tau_{\mathrm{apb}}}
\newcommand{\Ttimeout}{T_{\mathrm{timeout}}}
\newcommand{\Drescrow}{\Delta_{\mathrm{escrow}}}
\newcommand{\Dresume}{\Delta_{\mathrm{resume}}}
\newcommand{\Queue}{\mathcal{Q}}

% --- Title ---
\title{\textbf{Non-Blocking Governance:}\\[4pt]
       \large Escrow-Based Asynchronous Authorization for\\
       Human-Mediated Agent Oversight}

\author{Marcelo Fernandez\\
        \small TraslaIA\\
        \small \texttt{info@traslaia.com}}
\date{}

% =============================================================================
\begin{document}
\maketitle

% --- Abstract -----------------------------------------------------------------
\begin{abstract}
\noindent
The P9 MCP Governance Proxy establishes that safety can be enforced
protocol-transparently for any MCP-compatible agent---but at a cost:
when a tool call triggers a governance \HALT{}, the suspended call
blocks until a human returns a signed \emph{Accountability Proof Block}
(APB). In scenarios where human response is delayed or asynchronous,
this blocking model sacrifices liveness for safety.

We resolve this tension by introducing \emph{escrow-based non-blocking
governance}. When a risky tool call triggers a \HALT{}, its state is
serialised into a persistent escrow entry and deposited in a priority
queue rather than blocking the agent pipeline. The agent may continue
unrelated work while the suspended call remains in escrow. When the
APB arrives---in the same session or a different one---the escrow entry
is retrieved and the suspended tool call is resumed from the exact
preserved state.

We formalise three results. \emph{T10.1~(Non-Blocking Soundness)}
establishes that escrow preserves the core safety invariant: no risky
tool call $a \in A_{\mathrm{risky}}$ executes without a valid APB
satisfying predicates V1--V6, even while the agent continues other
tasks. \emph{T10.2~(Timeout Consistency)} establishes that when no
APB arrives within $\Ttimeout$, the fallback decision on the suspended
call is locally equivalent to an explicit governance decision for that
call. \emph{T10.3~(Escrow Liveness)} is the first liveness theorem in
the series: if the human signs within $\tsign \le \thalt + \Ttimeout$,
the suspended tool call resumes from the exact state
$\Eescrow(\thalt)$ with bounded latency~$\Dresume$.

Six experiments validate the construction.
\emph{E0}~(throughput, $N=100$ agents): P10 achieves
$10.7\times$--$81\times$ speedup over P9's blocking model as halt
rates rise from 10\% to 80\%, while unrelated safe-task throughput
remains unaffected by concurrent HALTs.
\emph{E1}~(overhead, $10{,}000$ cycles):
P95~$\Drescrow = 0.0064\,\si{\milli\second}$, $156\times$ below the
$1\,\si{\milli\second}$ gate.
\emph{E2}~(timeout semantics, $200\times3$ entries):
$0$ unauthorised executions across all three fallback modes.
\emph{E3}~(persistence, $N=100$, proxy restart simulation):
$100\%$ correct resume and $100\%$ V5 replay defence after session
boundary.
\emph{E4}~(concurrency, $N \in \{1,4,16,64\}$): $0$ exceptions,
$100\%$ APB validity, $0$ cross-agent contamination.
\emph{E5}~(adversarial): A1~(replay), A2~(post-timeout APB),
A3a~(policy forgery) all blocked at $0\%$ success; A3b~(queue-file
tampering) documents that the escrow store requires HMAC integrity
protection---an implementation gap addressed in P11.
\end{abstract}

% =============================================================================
\section{Introduction}
\label{sec:intro}

Human-mediated oversight of LLM agents faces a fundamental operational
constraint: human responses are asynchronous. A principal reviewing a
governance event---inspecting drift evidence, deliberating on risk,
consulting a policy document---takes seconds to minutes. An agent
waiting synchronously for that response is blocked: it cannot advance
any of its tasks, process other tool calls, or serve other users.

Paper~9 of this series~\citep{fernandez2026mcpgov} placed governance at
the MCP protocol layer, achieving zero-modification deployment and
transparent safety enforcement. But P9's architecture preserved P8's
blocking semantics~\citep{fernandez2026ibg}: \texttt{p9/apbRequired}
halts the agent pipeline until a signed APB is returned. As P9 explicitly
noted: \emph{``APBRequired prioritises safety over liveness. Paper~10 will
formalise timeout semantics, escrow-state serialisation, and
deferred-approval queues so that governance events can be resolved
asynchronously without halting the agent pipeline.''}

This paper fulfils that commitment.

\paragraph{Problem.}
Formal statement of the blocking limitation:

\begin{quote}
\emph{In the P9 model, a governance \HALT{} at time $\thalt$ blocks
progress on the affected tool call until a signed APB is returned,
reducing effective throughput on that call to zero during
$[\thalt,\, \thalt + \tau_{\mathrm{APB}}]$, where $\tau_{\mathrm{APB}}$
is the human response latency. As $\tau_{\mathrm{APB}} \to \infty$
(human offline), liveness of the governed call is lost entirely.}
\end{quote}

\paragraph{Contribution.}
We introduce \emph{escrow-based non-blocking governance}, a design
in which:
\begin{enumerate}[leftmargin=*, itemsep=2pt]
  \item A \HALT{} deposits the tool-call state into a persistent escrow
        store~$\Queue$ and immediately returns control to the agent.
  \item The agent may continue unrelated tasks while the suspended call
        remains in escrow; liveness for the suspended call is bounded by
        the TTL parameter~$\Ttimeout$, not by human response time.
  \item When the APB arrives (same session or a different one), the
        escrow entry is retrieved and the tool call is resumed with
        full cryptographic re-verification (V1--V6).
  \item If the APB does not arrive within~$\Ttimeout$, one of three
        configurable fallback modes fires: \DENY{} (safe default),
        \ADMIT{} (risk-accepted), or \ESCALATE{} (backup principal).
\end{enumerate}

The key invariant is that \emph{safety is not relaxed}: no risky
action executes without an explicit governance decision---either a valid
APB (V1--V6) or a fallback that itself constitutes a governance
decision. What changes is that this decision is no longer required to
be \emph{synchronous}.

\paragraph{Relationship to the series.}

\begin{center}
\small
\begin{tabular}{lll}
\toprule
Paper & What it formalises & Residual gap \\
\midrule
P8 & Identity-bound APB, V1--V5 & Blocking resolution \\
P9 & Protocol-layer proxy, T9.1--T9.3 & Liveness under async \\
\textbf{P10} & \textbf{Escrow, V6, T10.1--T10.3} & \textbf{Escrow integrity (A3b)} \\
P11 & Escrow-store HMAC, multi-hop formal & (future) \\
\bottomrule
\end{tabular}
\end{center}

% =============================================================================
\section{Background}
\label{sec:background}

\subsection{APB and Verification Predicates (P8)}

An \emph{Accountability Proof Block} is a triple
$\APB = (\Es, \Dh, \sigh)$ where $\Es$ is a system-constructed
evidence block, $\Dh$ is a human decision block, and $\sigh$ is an
ed25519 signature over $\mathrm{canonical}(\Es) \,||\,
\mathrm{canonical}(\Dh)$~\citep{fernandez2026ibg}. Five predicates
V1--V5 must hold for an APB to be valid:

\begin{itemize}[itemsep=1pt]
  \item \textbf{V1} Signature: $\sigh$ verifies under $\pki$ of the named principal.
  \item \textbf{V2} Registration: $\Hi \in \mathcal{P}$ (principal registry).
  \item \textbf{V3} Activity: $\Hi$ was active at~$t_e$ (not revoked).
  \item \textbf{V4} Freshness: $|t_{\mathrm{now}} - t_e| \le T_{\mathrm{max}}$ (replay window).
  \item \textbf{V5} Uniqueness: \texttt{event\_id} not previously seen (duplicate defence).
\end{itemize}

\subsection{MCP Governance Proxy (P9)}

$\Proxy$ intercepts every tool call between agent client and server.
On the non-\HALT{} path, $\Proxy$ is observationally equivalent to
direct pass-through (T9.1). On the \HALT{} path, $\Proxy$ emits
\texttt{p9/apbRequired} and blocks until \texttt{p9/apbResponse}
arrives~\citep{fernandez2026mcpgov}. The blocking latency is bounded by
$\Dnet + \Dverify + \Dpolicy$ (T9.2), but the human response time
$\tau_{\mathrm{APB}}$ is unbounded.

\subsection{The Liveness Gap}

Let $r$ be the fraction of tool calls that trigger a \HALT{} and
$T_{\mathrm{task}}$ the mean time per task. Under P9's blocking model,
agent throughput degrades as:
\[
  \Theta_{\mathrm{P9}}(r)
  = \frac{1}{T_{\mathrm{task}} + r \cdot \tau_{\mathrm{APB}}}.
\]
For $r = 0.5$ and $\tau_{\mathrm{APB}} = 2\,\si{s}$ (a modest human
response time), throughput drops by $99\%$ relative to $r=0$.
When the human is offline ($\tau_{\mathrm{APB}} \to \infty$),
$\Theta_{\mathrm{P9}} \to 0$ and liveness is lost entirely.

% =============================================================================
\section{System Model}
\label{sec:model}

\subsection{Escrow State}
\label{sec:escrow}

\begin{definition}[Escrow Entry]
An \emph{escrow entry} is a tuple
\[
  \Eescrow = (\mathit{tool\_call},\, \mathit{args},\, \mathit{context},\,
               \Dhat,\, \thalt)
\]
capturing the complete call state at the moment of \HALT{}. Fields:
\begin{itemize}[itemsep=1pt]
  \item $\mathit{tool\_call}$: the MCP tool name.
  \item $\mathit{args}$: the original call arguments (verbatim).
  \item $\mathit{context}$: session metadata (session identifier,
        trace hash, $\Azero$ hash).
  \item $\Dhat$: the IML drift estimator at~$\thalt$ (used for queue priority).
  \item $\thalt$: ISO~8601 UTC timestamp of the \HALT{} event.
\end{itemize}
Each entry is assigned an opaque \texttt{escrow\_id} (UUID4) on creation.
\end{definition}

Entries are serialised to JSON and stored in an \emph{EscrowStore}.
The serialisation is the identity function over field values.
In the current implementation, integrity protection is treated as a
future requirement rather than a protocol assumption
(cf.~\S\ref{sec:discussion} for the A3b implications).

\subsection{Priority Queue}
\label{sec:queue}

\begin{definition}[APB Queue]
The \emph{APB Queue} $\Queue$ is a priority store over escrow entries
ordered by $(-\Dhat,\, \thalt)$: highest drift first, then oldest
\HALT{} first (FIFO within equal drift). $\Queue$ supports:
\begin{itemize}[itemsep=1pt]
  \item $\Queue.\mathtt{enqueue}(e, \mathit{event\_id})$: add entry; raise
        \texttt{DuplicateEventError} if \texttt{event\_id} already pending.
  \item $\Queue.\mathtt{resolve}(id, \APB, \mathcal{P}, \pi)$:
        verify V1--V6; return $(d, e)$ where $d \in \{\RESUME, \DENY, \ESCALATE\}$.
  \item $\Queue.\mathtt{apply\_timeouts}(\pi, t)$: sweep expired entries;
        return $[(e_i, f_i)]$ where $f_i$ is the fallback mode.
  \item $\Queue.\mathtt{save}(p)$ / $\Queue.\mathtt{load}(p)$:
        atomic JSON persistence (multi-session support).
\end{itemize}
\end{definition}

\subsection{Timeout Policy and V6 Predicate}
\label{sec:timeout}

\begin{definition}[Timeout Policy]
A \emph{timeout policy} $\pi = (\Ttimeout, f)$ where:
\begin{itemize}[itemsep=1pt]
  \item $\Ttimeout \in (0, \infty]$: maximum seconds from $\thalt$ to
        APB arrival.
  \item $f \in \{\DENY, \ADMIT, \ESCALATE\}$: fallback action if
        $\Ttimeout$ expires without an APB.
\end{itemize}
\end{definition}

\begin{definition}[V6 Predicate --- Timeout Enforcement]
\label{def:v6}
Let $\tapb = \Es.t_e$ be the APB signing timestamp and
$\pi = (\Ttimeout, f)$ the active policy. Then:
\[
  V6(\APB, e, \pi)
  \;\triangleq\;
  \tapb \;\le\; \thalt + \Ttimeout.
\]
An APB passing V1--V5 but failing V6 is rejected with
\texttt{TimeoutExpiredError}; the entry remains pending for fallback
processing.
\end{definition}

V6 stacks on top of V5 in the verifier: an APB must satisfy all six
predicates for $\Queue.\mathtt{resolve}$ to succeed.
V4 defends against clock-based replay; V5 defends against
within-window duplicate submissions; V6 defends against late APBs
that would violate the $\Ttimeout$ commitment made to the agent.

\subsection{Fallback Modes}
\label{sec:fallback}

The three fallback modes represent distinct governance postures:

\begin{description}[itemsep=2pt]
  \item[\DENY{}] \emph{(default)} The tool call is dropped; no execution
    occurs. This preserves the ``no authority, no action'' invariant even
    in the absence of a human response. Appropriate for high-risk tools.
  \item[\ADMIT{}] The tool call is executed under an explicit policy
    decision that accepts the residual risk. It is appropriate only for
    low-risk tools where blocking would create unacceptable service
    degradation. \ADMIT{} is a governance posture, not a liveness
    mechanism.
  \item[\ESCALATE{}] The entry is forwarded to a backup principal or
    out-of-band resolution path without local execution. Appropriate for
    critical operations where both \DENY{} and \ADMIT{} are
    unsatisfactory.
\end{description}

% =============================================================================
\section{Non-Blocking Governance Protocol}
\label{sec:protocol}

\subsection{Session A: HALT and Escrow}

\begin{algorithm}[H]
\caption{Escrow enqueue on \HALT{} (replaces P9 blocking wait)}
\label{alg:enqueue}
\begin{algorithmic}[1]
\Require tool call $(a, \mathit{args})$, IML state $\Dhat$, queue $\Queue$, policy $\pi$
\State $t_{\mathrm{now}} \gets \mathtt{utc\_now}()$
\State $\Es \gets \mathtt{construct\_evidence}(\Azero, \Dhat, \mathit{trace}, t_{\mathrm{now}})$
\State $e \gets \Eescrow(a, \mathit{args}, \mathit{ctx}, \Dhat, t_{\mathrm{now}})$
\State $\mathit{eid} \gets \Queue.\mathtt{enqueue}(e,\; \Es.\mathtt{event\_id})$
\Comment{$O(1)$ --- returns immediately}
\State \textbf{emit} \texttt{apbRequired}$(a, \mathit{args}, \mathit{eid},\, \mathrm{summary}(\Es))$
\State \Return \textbf{control to agent}
\Comment{agent continues other tasks}
\end{algorithmic}
\end{algorithm}

\subsection{Session B: APB Arrival and Resume}

\begin{algorithm}[H]
\caption{Escrow resolution on APB arrival}
\label{alg:resolve}
\begin{algorithmic}[1]
\Require APB response $(\mathit{eid}, \APB)$, queue $\Queue$, registry $\mathcal{P}$, policy $\pi$
\State $e \gets \Queue.\mathtt{get}(\mathit{eid})$
\Comment{raises KeyError if not found}
\State $r \gets \mathtt{verify\_apb}(\APB, \mathcal{P})$
\Comment{V1--V5}
\If{$\neg r.\mathtt{is\_valid}$}
  \State \Return \texttt{REJECTED}, $r.\mathtt{result}$
\EndIf
\If{$\neg V6(\APB, e, \pi)$}
  \Comment{$\tapb > \thalt + \Ttimeout$}
  \State \Return \texttt{TIMEOUT\_EXPIRED}
\EndIf
\State $(d, e) \gets \Queue.\mathtt{resolve}(\mathit{eid}, \APB, \mathcal{P}, \pi)$
\State \textbf{resume} tool call $(e.\mathit{tool\_call}, e.\mathit{args})$
\Comment{from exact state $\Eescrow(\thalt)$}
\State \Return \texttt{RESOLVED}, $d$
\end{algorithmic}
\end{algorithm}

\subsection{Timeout Sweep}

At configurable intervals the queue runs a timeout sweep:
for each entry $e$ with $t_{\mathrm{now}} > \thalt + \Ttimeout$,
$\Queue.\mathtt{apply\_timeouts}(\pi, t_{\mathrm{now}})$ returns
$(e, f)$ where $f = \pi.f$. The caller is responsible for acting
on the fallback: discarding the call (\DENY), executing it (\ADMIT),
or forwarding it (\ESCALATE).

% =============================================================================
\section{Formal Results}
\label{sec:theory}

We state and prove three theorems. Proofs rely on the P8 APB
construction and P9 proxy invariants; references to earlier results
are cited explicitly.

\subsection{T10.1 --- Non-Blocking Soundness}

\begin{theorem}[Non-Blocking Soundness]
\label{thm:soundness}
Under the escrow protocol (Algorithms~\ref{alg:enqueue}--\ref{alg:resolve}),
for every risky tool call $a \in A_{\mathrm{risky}}$:
\[
  a \text{ executes}
  \;\Longrightarrow\;
  \exists\,\APB : V1(\APB) \land \cdots \land V6(\APB, e, \pi).
\]
Equivalently, no risky tool call executes without a valid APB satisfying
all six predicates or an explicit policy-authorised fallback, even when
the agent continues other tasks during the escrow period.
Furthermore, the protocol guarantees \emph{at-most-once} execution
semantics: each escrow entry is resumed at most once, enforced by the
combination of V5 (event-id uniqueness) and queue removal on
resolution.
\end{theorem}

\begin{proof}
By case analysis on the execution path.

\emph{Case 1: APB arrives before $\Ttimeout$.}
Algorithm~\ref{alg:resolve} runs V1--V5 (line~2) and V6 (line~5)
before executing the tool call (line~9). If any predicate fails, the
function returns early at lines 3 or 6 without executing $a$. The
tool call executes only if all six checks pass ($r.\mathtt{is\_valid}
\land V6$). This establishes the implication.

\emph{Case 2: $\Ttimeout$ expires, fallback $= \DENY$.}
$\Queue.\mathtt{apply\_timeouts}$ removes the entry and returns
$(e, \DENY)$. The tool call is not executed; the entry is permanently
removed from $\Queue$. No path remains by which $a$ can execute.

\emph{Case 3: $\Ttimeout$ expires, fallback $= \ADMIT$.}
The suspended tool call executes only because the configured policy
explicitly authorises it. This constitutes an explicit policy
authorisation rather than an unauthorised execution.

\emph{Case 4: $\Ttimeout$ expires, fallback $= \ESCALATE$.}
The entry is forwarded to a backup principal. No direct execution
occurs until that principal responds with a valid APB, which reduces
to Case~1.

In all cases, the suspended tool call does not execute without either
a valid APB satisfying V1--V6 or an explicit policy-authorised
fallback for that suspended call.
\qed
\end{proof}

\begin{remark}
T10.1 extends P9's Halt Invariant to the async setting. The key
difference: in P9, safety follows from blocking (no APB = no
execution = no progress). In P10, safety follows from the escrow
invariant: a tool call can leave escrow only through a resolved APB
(cases 1, 3, 4) or a \DENY{} fallback (case 2).
\end{remark}

\subsection{T10.2 --- Timeout Consistency}

\begin{theorem}[Timeout Consistency]
\label{thm:timeout}
Let $e = \Eescrow(\thalt)$ be an escrow entry and $\pi = (\Ttimeout, \DENY)$.
If no valid APB arrives within $\Ttimeout$, the observable state of the
\emph{suspended tool call, considered in isolation}, is equivalent at
timeout expiration to its state immediately post-\HALT{}:
\[
  \mathrm{state}_{a}(\thalt + \Ttimeout)
  \;\equiv_{\mathrm{obs}}\;
  \mathrm{state}_{a}(\thalt).
\]
Furthermore, the total latency to resume after a valid APB arrives at
$\tsign \le \thalt + \Ttimeout$ is bounded:
\[
  \Dresume \;\le\; \Drescrow + \Delta_{\mathrm{get}},
\]
where $\Drescrow$ is the serialisation cost and $\Delta_{\mathrm{get}}$
is the deserialisation and in-memory lookup cost.
\end{theorem}

\begin{proof}
\emph{Observational equivalence at DENY.}
When $\DENY$ fires, $\Queue.\mathtt{apply\_timeouts}$ removes the entry
from the pending queue without altering the suspended call state. The
suspended call state $(\mathit{tool\_call}, \mathit{args},
\mathit{context}, \Dhat)$ is unchanged from its value at enqueue (which
equals $\thalt$). Since no external side effect of the suspended tool
call has occurred (Theorem T10.1, Case~2), the observable state is
identical to what it was at $\thalt$.

\emph{Bounded resume latency.}
The entry $e$ is stored via $\Eescrow.\mathtt{to\_json}()$ (cost
$\Drescrow$) and retrieved via $\mathtt{EscrowStore.get}()$ (cost
$\Delta_{\mathrm{get}}$). Both operations are $O(1)$ in the size of
the entry (constant-width fields). By the E1 measurement
(\S\ref{sec:experiments}), the P95 round-trip cost is
$0.0064\,\si{\milli\second}$, confirming that $\Dresume$ is bounded
and negligible relative to LLM inference and human response times.
\qed
\end{proof}

\begin{remark}
T10.2 covers only the state of the suspended \emph{tool call}, not the
state of the agent as a whole. The agent may have made progress on
other tasks during $[\thalt, \thalt + \Ttimeout]$; T10.2 makes no
claim about that progress. This is intentional: the equivalence is
local to the suspended call, which is the correct scope for a
governance consistency theorem.
\end{remark}

\subsection{T10.3 --- Escrow Liveness}

\begin{theorem}[Escrow Liveness]
\label{thm:liveness}
Let $e = \Eescrow(\thalt)$ be an escrow entry, $\pi = (\Ttimeout, f)$
a timeout policy, and $\APB$ a valid APB (V1--V6) signed at
$\tsign \le \thalt + \Ttimeout$.  Then the suspended tool call resumes
within bounded latency $\Dresume$ after APB submission at $\tsign$:
\[
  \tsign \le \thalt + \Ttimeout
  \;\Longrightarrow\;
  \text{tool call resumes within } \tsign + \Dresume.
\]
\end{theorem}

\begin{proof}
By the V6 predicate (Definition~\ref{def:v6}):
$\tsign \le \thalt + \Ttimeout$ is precisely the condition $V6(\APB, e,
\pi) = \top$.  Combined with V1--V5 (which hold by assumption that
$\APB$ is valid), Algorithm~\ref{alg:resolve} reaches line~9 and
executes the tool call.

The execution happens within $\tsign + \Dresume$ because:
(a) the APB is submitted at $\tsign$;
(b) V1--V5 take $O(\Dverify)$, V6 is a single timestamp comparison
    $O(1)$;
(c) the entry is retrieved from $\Queue$ in $O(\Delta_{\mathrm{get}})$.

Thus the total delay from APB submission to tool call execution is
$\Dverify + \Delta_{\mathrm{get}} \le \Dresume$, which is empirically
bounded by $0.0064\,\si{\milli\second}$ P95 (E1,
\S\ref{sec:experiments}).

This is the first \emph{liveness} theorem in the series. It does not
remove the human-responsiveness assumption; it makes that assumption
explicit and checkable through V6 rather than leaving it as an
implicit prerequisite.
\qed
\end{proof}

\begin{remark}
T10.3 is conditional on \emph{human-responsiveness}: the bound
$\tsign \le \thalt + \Ttimeout$ must hold. This is the unavoidable
cost of human-in-the-loop governance. The contribution of T10.3 is
not to eliminate this dependency, but to show that it can be
characterised as a bounded, checkable condition---V6---rather than
an implicit assumption.
\end{remark}

% =============================================================================
\section{Implementation}
\label{sec:impl}

The escrow layer is implemented in Python~3.12 and comprises three
modules:

\begin{description}[itemsep=4pt]
  \item[\texttt{escrow/escrow\_store.py}]
    \texttt{EscrowEntry} is a \texttt{frozen dataclass} with six
    fields matching Definition~3.1. JSON serialisation uses stdlib
    \texttt{json} (no external dependencies). \texttt{EscrowStore}
    provides \texttt{put}/\texttt{get}/\texttt{remove}/\texttt{list\_entries}
    plus atomic file persistence via \texttt{save(path)} /
    \texttt{load(path)} (write-to-tmp then rename).

  \item[\texttt{escrow/timeout\_policy.py}]
    \texttt{FallbackMode} (\DENY/\ADMIT/\ESCALATE) is a \texttt{str}
    enum. \texttt{TimeoutPolicy} holds $(\Ttimeout, f)$ and exposes
    \texttt{is\_expired}(\textit{entry}, \textit{now}) and
    \texttt{apply\_fallback}(\textit{entry}).
    \texttt{v6\_check}(\textit{apb\_t\_e}, $\thalt$, $\Ttimeout$) is a
    standalone function imported by \texttt{APBQueue.resolve}.

  \item[\texttt{escrow/apb\_queue.py}]
    \texttt{APBQueue} wraps \texttt{EscrowStore} with priority ordering,
    event-id deduplication (two dictionaries: \textit{event\_to\_escrow}
    and \textit{escrow\_to\_event}), and a \texttt{\_resolved\_event\_ids}
    set for V5 replay defence that persists across sessions.
    \texttt{resolve} calls V1--V5 via the frozen P9 verifier, then V6
    via \texttt{v6\_check}. Three custom exceptions carry structured
    diagnostic data: \texttt{DuplicateEventError},
    \texttt{APBVerificationError}, \texttt{TimeoutExpiredError}.
\end{description}

The frozen P7--P9 stack (\texttt{stack/}, \texttt{proxy/},
\texttt{agent/}) is untouched; the escrow layer extends it with zero
modifications to existing files.

% =============================================================================
\section{Experiments}
\label{sec:experiments}

All experiments run on Python~3.12, Windows~11, Intel architecture.
Results are deterministic (fixed seed where applicable). Raw data
in \texttt{results/}.

\subsection{E0 --- Throughput Comparison (P9 vs.\ P10)}

\paragraph{Setup.}
$N=100$ agents, $M=200$ tasks per agent, $T_{\mathrm{task}}=10\,\si{\milli\second}$,
$\tau_{\mathrm{APB}}=2\,\si{s}$ (human response latency), halt rates
$r \in \{0, 0.1, 0.2, 0.4, 0.6, 0.8\}$. Logical-time simulation
(no wall-clock sleep) using measured $\Drescrow$ and $\Dresume$ P95
values as escrow overhead.

\paragraph{Two metrics (measured separately).}
\begin{itemize}[itemsep=2pt]
  \item \textbf{Metric 1} ($\Dresume$): latency from APB arrival to
    task resumption; measured empirically in E1 ($\Delta_{\mathrm{resume}}$
    P95 $= 0.0001\,\si{\milli\second}$).
  \item \textbf{Metric 2} (continuity): throughput of non-blocked tasks
    during the escrow window. In P9, all tasks block ($\Theta_{\mathrm{safe}}
    = 0$). In P10, safe tasks are unaffected ($\Theta_{\mathrm{safe}} =
    1/T_{\mathrm{task}} = 100\,\text{tps}$).
\end{itemize}

\begin{table}[h]
\centering
\caption{E0 --- Throughput comparison (tasks per second, mean over $N=100$ agents).
P10~total includes both safe tasks and resumed escrow tasks.
P10~safe is the throughput of non-blocked tasks \emph{during} the escrow window.}
\label{tab:e0}
\begin{tabular}{rrrrr}
\toprule
Halt rate & P9 (tps) & P10 total (tps) & P10 safe (tps) & Speedup \\
\midrule
0.0 & 100.00 & 100.00 & 100.00 & $1.00\times$ \\
0.1 &   4.76 &  51.06 & 100.00 & $10.7\times$ \\
0.2 &   2.49 &  50.53 & 100.00 & $20.3\times$ \\
0.4 &   1.25 &  50.21 & 100.00 & $40.0\times$ \\
0.6 &   0.84 &  50.07 & 100.00 & $59.8\times$ \\
0.8 &   0.62 &  50.04 & 100.00 & $81.0\times$ \\
\bottomrule
\end{tabular}
\end{table}

P10 total throughput is bounded below by $N_{\mathrm{safe}}/t_{\mathrm{agent}}$,
which approaches $50\,\text{tps}$ as $r \to 1$ (half the tasks are safe,
none block). P9 throughput degrades as
$\Theta_{\mathrm{P9}} \approx 1/(r \cdot \tau_{\mathrm{APB}})$ for large~$r$.

\paragraph{Gate.} P10 throughput $\ge$ P9 throughput at all $r > 0$:
\textbf{PASSED}.

\subsection{E1 --- Escrow Overhead}

\paragraph{Setup.} $10{,}000$ cycles per benchmark operation. Table~\ref{tab:e1}.

\begin{table}[h]
\centering
\caption{E1 --- Escrow operation latencies ($\si{\milli\second}$, Python~3.12).}
\label{tab:e1}
\begin{tabular}{lrrrr}
\toprule
Operation & Mean & P50 & P95 & P99 \\
\midrule
Roundtrip (ser.\ + deser.) & 0.0059 & 0.0056 & \textbf{0.0064} & 0.0095 \\
Serialize only              & 0.0029 & 0.0028 & 0.0031 & 0.0036 \\
Deserialize only            & 0.0029 & 0.0028 & 0.0031 & 0.0036 \\
Put (in-memory)             & 0.0002 & 0.0001 & 0.0003 & 0.0005 \\
\midrule
File save (1 entry)         & 0.3572 & 0.3388 & 0.4562 & 0.6404 \\
File save (100 entries)     & 1.1253 & 1.0871 & 1.3095 & 1.4017 \\
File load (1 entry)         & 0.0833 & 0.0462 & 0.0806 & 0.1750 \\
File load (100 entries)     & 0.3242 & 0.2835 & 0.3611 & 0.4599 \\
\bottomrule
\end{tabular}
\end{table}

P95 roundtrip $= 0.0064\,\si{\milli\second}$, which is $156\times$ below
the $1\,\si{\milli\second}$ gate. File save/load costs are dominated
by OS I/O and are negligible relative to LLM inference
($\sim$seconds).

\paragraph{Gate.} P95~$\Drescrow < 1\,\si{\milli\second}$:
\textbf{PASSED} ($0.0064\,\si{\milli\second}$).

\subsection{E2 --- Timeout Semantics}

\paragraph{Setup.}
$N=200$ expired entries per fallback mode
($\Ttimeout=60\,\si{s}$, entries created $120\,\si{s}$ before
sweep time). Mixed-queue trial: $50$ expired $+$ $50$ active entries.

All 200 entries swept per mode; 0 unauthorized executions across
\DENY, \ADMIT, \ESCALATE. Active entries in the mixed trial were
not swept ($50/50$ preserved). An \emph{``unauthorised execution''}
is defined as any entry that exits the queue without receiving a
governance decision (valid APB or explicit fallback); the count was~$0$
in all trials.

\paragraph{Gate.} $0$ unauthorized executions: \textbf{PASSED}.

\subsection{E3 --- Queue Persistence}

\paragraph{Setup.}
$N=100$ entries; full-restart trial (all entries enqueued in session~A,
saved, reloaded in session~B, resolved with fresh APBs); partial-restart
trial (half resolved pre-restart, remainder post-restart with V5 replay
test).

Full-restart: $100/100$ field-parity checks, $100/100$ correct
governance decisions, $0$ entries remaining after resolution.
Partial-restart: $50/50$ active entries present after reload;
$50/50$ pre-restart APBs blocked by V5
(\texttt{DUPLICATE\_EVENT\_ID}) after session~B reloads the
\texttt{\_resolved\_event\_ids} set.

\paragraph{Gate.} $100\%$ correct resume post-restart: \textbf{PASSED}.

\subsection{E4 --- Concurrency}

\paragraph{Setup.}
$N \in \{1,4,16,64\}$ agents, each with $M=20$ entries, all
running concurrently (\texttt{ThreadPoolExecutor}). Each agent owns
an independent \texttt{APBQueue} (architecture reflects one queue per
proxy instance, not shared state).

$0$ exceptions across all configurations; $100\%$ APB validity;
$0$ cross-agent entry contamination. At $N=64$: $1{,}280$ total
resolutions, all correct.

\paragraph{Gate.} $0$ exceptions, $100\%$ APB validity:
\textbf{PASSED}.

\subsection{E5 --- Adversarial Async}

\paragraph{Setup.} $N=100$ attacks per vector.

\begin{description}[itemsep=3pt]
  \item[A1 --- Replay event\_id.] After legitimate resolution, attacker
    re-submits the same APB to a fresh entry. Blocked by V5
    (\texttt{DUPLICATE\_EVENT\_ID}); $100/100$ blocked.

  \item[A2 --- APB post-timeout.] APB signed at
    $\tapb = t_{\mathrm{halt}} + 90\,\si{\min}$ against a
    $\Ttimeout=60\,\si{\min}$ policy.
    Blocked by V6 (\texttt{TimeoutExpiredError}); $100/100$ blocked.

  \item[A3a --- Policy forgery.] Attacker supplies a lenient
    $\Ttimeout = 10^{12}\,\si{s}$; server enforces strict
    $\Ttimeout = 3{,}600\,\si{s}$. Since the server controls
    the policy parameter in \texttt{resolve()}, the attack is
    impossible architecturally; $100/100$ blocked.

  \item[A3b --- Queue file tampering.] Attacker modifies the
    \texttt{escrow.json} file on disk, extending \texttt{t\_halt}
    to make a late APB appear within the window.
    \textbf{Finding:} tampering succeeded in $100/100$ trials.
    This is documented as an open security requirement: the escrow
    store requires HMAC integrity protection (see~\S\ref{sec:discussion}).
\end{description}

\paragraph{Gate.} $0\%$ adversary success (A1/A2/A3a):
\textbf{PASSED}. A3b: \textbf{documented finding}.

% =============================================================================
\section{Discussion}
\label{sec:discussion}

\subsection{The A3b Finding: Escrow Integrity}

The A3b result is not a defect in the escrow protocol; it is an
\emph{implementation gap} in the current persistence layer. V6 is
correct only if the stored escrow metadata is authentic. The
experiment shows that plaintext JSON storage does not provide that
guarantee. This is not a weakness of the escrow model itself, but of
the current persistence layer.

The appropriate mitigation is to make the escrow store
\emph{tamper-evident}, analogous to P8's \texttt{APBLog}: HMAC-chain
persisted entries with a server-held key. An attacker without the key
cannot modify \texttt{t\_halt}---or any other escrow field---without
invalidating the chain. Formally:

\begin{quote}
\emph{Requirement R10.1:} The \texttt{EscrowStore} must maintain an
integrity-protected chain over persisted entries such that any
post-write modification is detected with probability
$1 - 2^{-\lambda}$, where $\lambda$ is the integrity parameter.
The protocol-level guarantees of V6 assume authentic escrow metadata;
R10.1 is the implementation requirement that backs that assumption.
\end{quote}

This requirement will be formalised and implemented in P11 alongside
the multi-hop formal treatment. We include A3b precisely to make this
gap explicit rather than implicit.

\subsection{Non-Blocking vs.\ Weaker Guarantees}

A na\"ive alternative to escrow is a simple timeout: if no APB arrives
within $\Ttimeout$, deny the call and continue. This achieves liveness
but loses T10.3's resume guarantee: the tool call state is not preserved,
and the agent must re-derive it on retry (which may not be possible if
the original context has changed). Escrow is strictly stronger: it
preserves the exact call state and enables \emph{replay from
$\Eescrow(\thalt)$} rather than retry.

\subsection{Relationship to CAP-Like Tradeoffs}

An intuition from distributed systems is useful here, though we do not
formalise it as a CAP theorem (the definitions are not directly
applicable):

\begin{itemize}[itemsep=2pt]
  \item P9 prioritises \emph{safety} (C in the analogy) over
    \emph{liveness} (A): blocking until APB arrives.
  \item P10 achieves \emph{bounded liveness} without relaxing safety:
    the escrow invariant ensures safety holds at all times; T10.3
    bounds liveness under the human-responsiveness assumption.
\end{itemize}

The key insight is that safety and liveness are not fundamentally
opposed in human-mediated governance---they can coexist if the
human response is bounded. V6 is the formal expression of that bound.

\subsection{Scope of T10.2}

T10.2's observational equivalence is local to the suspended
\emph{tool call}, not the agent as a whole. During $[\thalt,
\thalt + \Ttimeout]$, the agent may complete other tasks, producing
observable side effects. T10.2 does not claim those are absent; it
claims only that the \emph{state of the suspended call} is preserved.
This is the correct scope: governance is per-call, not per-session.

% =============================================================================
\section{Related Work}
\label{sec:related}

\paragraph{LLM agent governance.}
The Agent Governance Series~\citep{fernandez2026a,fernandez2026acp,
fernandez2026iml,fernandez2026govstr,fernandez2026ram,fernandez2026opram,
fernandez2026applied,fernandez2026ibg,fernandez2026mcpgov} provides the
foundational stack that P10 extends. The blocking limitation of P9 is
the direct motivation for this work.

\paragraph{Runtime verification.}
Leucker and Schallhart~\citep{leucker2009brief} survey runtime
verification; our IML monitor (P2) fits this framework.
Basin et al.~\citep{basin2012monitoring} study monitoring for
usage-control policies in distributed systems, a close analogue of
our principal-based governance.

\paragraph{Asynchronous authorization.}
OAuth~2.0 and OpenID Connect \citep{hardt2012oauth} support
asynchronous token issuance but do not formalise liveness bounds
under delayed responses. Our V6 predicate and T10.3 fill this gap
for the agent governance setting.

\paragraph{Persistence and recovery.}
The HMAC-chained log of P8~\citep{fernandez2026ibg} and the write-ahead
log in database systems~\citep{mohan1992aries} both use
tamper-evident persistence; the A3b finding shows P10 needs to apply
the same principle to the escrow store.

% =============================================================================
\section{Conclusion}
\label{sec:conclusion}

We have introduced escrow-based non-blocking governance, a design that
preserves the P8--P9 safety guarantee---no risky tool call executes
without a valid APB satisfying V1--V6---while adding bounded liveness
under asynchronous human resolution. The three theorems establish:

\begin{itemize}[itemsep=2pt]
  \item \textbf{T10.1} (Soundness): escrow does not relax safety;
    each escrow entry is resumed at most once.
  \item \textbf{T10.2} (Consistency): timeout fallback is a local
    governance decision on the suspended tool call; resume latency
    is bounded.
  \item \textbf{T10.3} (Liveness): the first liveness theorem in the
    series; when the human responds within $\Ttimeout$, the suspended
    call resumes with bounded delay
    $\Dresume = 0.0064\,\si{\milli\second}$ P95.
\end{itemize}

The experiments support the construction and isolate one remaining
implementation requirement: escrow persistence must be
integrity-protected before V6 can be considered tamper-resistant in
production. The A3b finding makes this gap explicit and provides a
clear scope for P11.

P10 closes the liveness gap left open by P9 and establishes the formal
foundation for the next Phase~II papers---P11 (escrow integrity and
multi-hop formal treatment) and P12 (distributed multi-organisation
governance)---while leaving escrow-store integrity as an explicit
future requirement rather than a hidden assumption.

% =============================================================================
\bibliographystyle{plainnat}
\bibliography{references}

\end{document}
